\section{The final pairing}\label{sec:conclusion}

\begin{proof}[Proof of Theorem~\ref{thm:main}]
Let \(\omega\) be a smooth closed form taming \(J\).
Suppose for contradiction that there is no compatible
symplectic form, and take \(P,Q,T,\lambda,\mu\) from
Proposition~\ref{prop:separation}. Propositions
\ref{prop:growth} and \ref{prop:energy} verify all the
hypotheses of Theorem~\ref{thm:splitting}.
Consequently the current \(P_E=\ind_{\{\theta>0\}}P\)
is closed. Set
\[
P'=\ind_{X\setminus E}P,\qquad
T'=P'+Q=T-P_E.
\]
Then \(T'\) is closed. The current \(P'\) is represented
by the restriction \(\lambda'\le\lambda\) to base points
outside \(E\), with projection
\(\mu'=\ind_{X\setminus E}\mu\).
We pair \(T\) with this residual current rather than with itself:
placing the second positive factor on zero-density centers makes
the previously bounded negative error in the pairing tend to zero.

We first show that the negative term in
\eqref{eq:positive-pairing} tends to zero for this
choice of \(P'\). For each \(y\notin E\),
\begin{equation}\label{eq:zero-density-tail}
r^{-2}\int_X(1+d_g(x,y)/r)^{-6}\dd\mu(x)
\longrightarrow0.
\end{equation}
On the inner ball and on every fixed dyadic shell this
follows from \(s^{-2}\mu(B_s(y))\to0\).
The quadratic mass bound dominates the shell contributions
by \(C2^{-4j}\), a summable sequence, giving the full
limit. The uniform bound \eqref{eq:shell} then permits
dominated convergence against \(\mu'\).
Discarding the nonnegative angular term in
\eqref{eq:positive-pairing} yields
\begin{equation}\label{eq:liminf}
\liminf_{r\downarrow0}\ip{S_rP}{*S_rP'}_{L^2}\ge0.
\end{equation}

Apply Lemma~\ref{lem:zero-pairing} to \(T'\).
Its regularizations are in \(L^2\), since \(P'\) has
measure coefficients and \(Q\in L^2\).
Expanding gives
\[
0=\ip{S_rP}{*S_rP'}
+\ip{S_rP}{*S_rQ}
+\ip{S_rQ}{*S_rP'}
+\|S_rQ\|_2^2.
\]
Both mixed terms tend to zero by
\eqref{eq:cross-vanish}, applied once to \(P\) and once
to \(P'\). Since \(S_rQ\to Q\) strongly in \(L^2\),
\eqref{eq:liminf} implies \(0\ge\|Q\|_2^2\).
Therefore \(Q=0\).

It follows that \(P=T\) annihilates all smooth closed
forms. In particular \(P(\omega)=0\). On the other
hand, \(\omega_x(\ell)>0\) on the compact bundle
\(\mathscr C\); it has a positive minimum there.
Equation~\eqref{eq:positive-representation} and
\(\lambda(\mathscr C)=1\) give \(P(\omega)>0\),
a contradiction.

The positive cone and the subspace of closed invariant
smooth forms therefore intersect. Any form in their
intersection is the required smooth compatible
symplectic form for \(J\).
\end{proof}
